%%% FIELD ipw-upper-statement
Uniformly over \(P\in\mathfrak P^{\mathrm{pred}}_{n,m}(r,\mu,K,q,B,\lambda)\), define
\[
 a^{\mathrm{loose}}_{nm}=
 \sqrt{\frac n{\mu r}}\left[1\wedge
 \frac B{\lambda_{nm}}\left\{
 \sqrt{\frac{(n\vee m)\log(e(n+m))}{q}}+
 \frac{\log(e(n+m))}{q}\right\}\right].
\]
Then recorded-propensity IPW--SVD satisfies
\(\mathcal L(\widehat U^{\mathrm I},U)=O_P(a^{\mathrm{loose}}_{nm})\).  This maximum-row assertion is pointwise in the target and contains no factor depending on \(|\mathcal G_n|\).  In general the displayed IPW--SVD conclusion cannot replace \(a^{\mathrm{loose}}_{nm}\) by the lower/reference scale \(b_{nm}\), as witnessed by Proposition \(\mathrm{prop:adaptive\mbox{-}endogenous\mbox{-}target\mbox{-}obstruction}\).  That witness concerns this estimator only and does not assert that every observed-data estimator has risk larger than order \(b_{nm}\).
%%% FIELD ipw-upper-proof
Proposition \(\mathrm{prop:ipw\mbox{-}expected\mbox{-}risk\mbox{-}upper}\) gives the stronger uniform first-moment statement
\[
 \sup_{P\in\mathfrak P^{\mathrm{pred}}_{n,m}}
 E_P\!\left[\mathcal L(\widehat U^{\mathrm I},U)\right]
 \le C a^{\mathrm{loose}}_{nm}.
\]
The radius is strictly positive because \(B>0\), \(q>0\), \(\lambda_{nm}>0\), and \(n,m\ge2\).  Therefore Markov's inequality gives, for every \(t>0\),
\[
 \sup_{P\in\mathfrak P^{\mathrm{pred}}_{n,m}}
 P\!\left\{\mathcal L(\widehat U^{\mathrm I},U)>
 t a^{\mathrm{loose}}_{nm}\right\}\le \frac Ct.
\]
This is the asserted uniform \(O_P(a^{\mathrm{loose}}_{nm})\) bound.  The construction of \(a^{\mathrm{loose}}_{nm}\) and the proof of the moment bound use only the score matrix and contain no loading family, so no factor depending on \(|\mathcal G_n|\) occurs.

For the final assertion, Proposition \(\mathrm{prop:adaptive\mbox{-}endogenous\mbox{-}target\mbox{-}obstruction}\) supplies a sequence \(P_n\) in the same class and constants \(c_0,c_1>0\) such that \(b_{nn}\to0\) and
\[
 \liminf_{n\to\infty}P_n\!\left\{
 \mathcal L(\widehat U^{\mathrm I},U)>c_0\right\}\ge c_1.
\]
If the IPW--SVD loss were uniformly \(O_P(b_{nm})\), then for \(\varepsilon=c_1/2\) some fixed \(T<\infty\) would make the displayed probability at threshold \(Tb_{nn}\) at most \(\varepsilon\) eventually.  Since \(Tb_{nn}<c_0\) eventually, this contradicts the preceding lower probability.  The witness evaluates \(\widehat U^{\mathrm I}\) and hence proves no impossibility statement for a different estimator.
%%% FIELD expected-risk-statement
Put
\[
 a^{\mathrm{loose}}_{nm}=
 \sqrt{\frac n{\mu r}}\left[1\wedge
 \frac B{\lambda_{nm}}\left\{
 \sqrt{\frac{(n\vee m)\log(e(n+m))}{q}}+
 \frac{\log(e(n+m))}{q}\right\}\right].
\]
There is a universal constant \(C<\infty\) such that, for every \(n,m\ge2\),
\[
 \sup_{P\in\mathfrak P^{\mathrm{pred}}_{n,m}(r,\mu,K,q,B,\lambda)}
 E_P\!\left[\mathcal L(\widehat U^{\mathrm I},U)\right]
 \le C a^{\mathrm{loose}}_{nm}.
\]
Thus the loose IPW--SVD upper bound holds in the same expected normalized maximum-row loss used by the minimax lower benchmark.
%%% FIELD expected-risk-proof
Fix a law \(P\) in the predictable class and expose cells in the row-major order of \(\mathrm{def:global\mbox{-}filtration}\).  By Lemma \(\mathrm{lem:ipw\mbox{-}identification}\), each \(E_{ij}\) has conditional mean zero given \(\mathbb G_{o(i,j)-1}\).  Sequential randomization and bounded outcomes give
\[
 \begin{aligned}
 E_P\!\left[(Z^{\mathrm{ipw}}_{ij})^2
 \mid\mathbb G_{o(i,j)-1}\right]
 &=E_P\!\left[
 \frac{Y_{ij}(1)^2}{p_{ij}}+
 \frac{Y_{ij}(0)^2}{1-p_{ij}}
 \mathrel{\Big|}\mathbb G_{o(i,j)-1}\right]\\
 &\le \frac{2B^2}{q}.
 \end{aligned}
\]
The equality follows by first conditioning on the two potential outcomes and using sequential randomization; both divisions are valid by overlap.  Since \(M_{ij}=E_P[Z^{\mathrm{ipw}}_{ij}\mid\mathbb G_{o(i,j)-1}]\), conditional variance is no larger than the displayed conditional second moment.  Moreover \(|Z^{\mathrm{ipw}}_{ij}|\le B/q\), \(|M_{ij}|\le2B\), and \(q<1/2\), whence
\[
 |E_{ij}|\le \frac{2B}{q}.
\]

Introduce the self-adjoint dilation
\[
 \Delta_{ij}=
 \begin{pmatrix}
 0&E_{ij}e_ie_j^\top\\
 E_{ij}e_je_i^\top&0
 \end{pmatrix}.
\]
Its norm is \(|E_{ij}|\), and its square has diagonal blocks
\(E_{ij}^2e_ie_i^\top\) and \(E_{ij}^2e_je_j^\top\).  Therefore the predictable quadratic variation of the sum of the dilations obeys
\[
 \left\lVert\sum_{i,j}E_P[\Delta_{ij}^2
 \mid\mathbb G_{o(i,j)-1}]\right\rVert
 \le v_{nm}:=\frac{2B^2(n\vee m)}q.
\]
Lemma \(\mathrm{lem:tropp\mbox{-}matrix\mbox{-}freedman}\), applied to both signs of this \((n+m)\)-dimensional martingale, consequently gives
\[
 P\{\lVert E\rVert\ge t\}
 \le2(n+m)\exp\!\left\{-\frac{t^2/2}{v_{nm}+R_{nm}t/3}\right\},
 \qquad R_{nm}=\frac{2B}{q}.
\]
Indeed, the terminal dilation has operator norm \(\lVert E\rVert\).  Solving the quadratic inequality in this tail bound shows that, with \(L_{nm}=\log(2(n+m))\), for every \(x\ge0\),
\[
 P\!\left\{\lVert E\rVert>
 \sqrt{2v_{nm}(L_{nm}+x)}+
 \frac{2R_{nm}}3(L_{nm}+x)\right\}\le e^{-x}.
\]
Integration of this bound is legitimate because the increments are bounded.  Using
\(\sqrt{L_{nm}+x}\le\sqrt{L_{nm}}+\sqrt x\),
\(\int_0^\infty e^{-x}\,dx=1\), and
\(\int_0^\infty e^{-x}\sqrt x\,dx<\infty\), yields
\[
 E_P\lVert E\rVert
 \le C B\left\{
 \sqrt{\frac{(n\vee m)\log(e(n+m))}{q}}+
 \frac{\log(e(n+m))}{q}\right\}
 =:C\eta_{nm}.
\]
Every bound so far is uniform in \(P\).

It remains to pass from operator error to the declared loss without imposing predictability on the singular factors.  Let \(P_U=UU^\top\), and let \((\widehat U,\widehat\Sigma,\widehat V)\) be the leading rank-\(r\) singular triple of \(M+E\).  Exact rank gives \((I-P_U)M=0\), while the signal floor and the singular-value min--max formula imply, on \(\{\lVert E\rVert\le\lambda_{nm}/2\}\),
\[
 \sigma_r(\widehat\Sigma)
 \ge\sigma_r(M)-\lVert E\rVert
 \ge\lambda_{nm}-\lVert E\rVert>0.
\]
Projecting \((M+E)\widehat V=\widehat U\widehat\Sigma\) onto the orthogonal complement of \(U\) therefore gives
\[
 \lVert(I-P_U)\widehat U\rVert
 \le\frac{\lVert E\rVert}{\lambda_{nm}-\lVert E\rVert}
 \le\frac{2\lVert E\rVert}{\lambda_{nm}}.
\]
Diagonalizing the principal-angle cosines of \(U^\top\widehat U\) and using
\(2(1-\cos t)\le2\sin^2t\) supplies an \(R\in\mathbb O_r\) such that
\[
 \lVert\widehat UR-U\rVert
 \le\sqrt2\lVert(I-P_U)\widehat U\rVert.
\]
Off the event \(\{\lVert E\rVert\le\lambda_{nm}/2\}\), the left side is at most two.  Since \(\lVert Q\rVert_{2,\infty}\le\lVert Q\rVert\), both cases combine to give the deterministic inequality
\[
 \mathcal L(\widehat U^{\mathrm I},U)
 \le C\sqrt{\frac n{\mu r}}
 \left(1\wedge\frac{\lVert E\rVert}{\lambda_{nm}}\right).
\]
The fixed tie rule ensures that the leading left singular vectors of \(\mathcal P_r(M+E)\) are the chosen leading left singular vectors of \(M+E\).  Finally, concavity of \(x\mapsto1\wedge x\), Jensen's inequality, and the expectation bound above yield
\[
 \begin{aligned}
 E_P\mathcal L(\widehat U^{\mathrm I},U)
 &\le C\sqrt{\frac n{\mu r}}
 E_P\left[1\wedge\frac{\lVert E\rVert}{\lambda_{nm}}\right]\\
 &\le C\sqrt{\frac n{\mu r}}
 \left(1\wedge\frac{\eta_{nm}}{\lambda_{nm}}\right)
 \le C a^{\mathrm{loose}}_{nm}.
 \end{aligned}
\]
Taking the supremum over the declared class proves the result.
%%% FIELD minimax-bracket-statement
Fix \(r\), \(K\), and \(\mu>1\).  There are constants \(c_0,c_1,c,C>0\), depending only on these fixed quantities, \(q\), and the aspect bounds, such that, for all sufficiently large \(n,m\), whenever \(b_{nm}\le c_0\) and \(\lambda_{nm}\le c_1B\sqrt{nm}\),
\[
 c b_{nm}\le
 \inf_{\widetilde U\in\mathfrak E_{n,m}}
 \sup_{P\in\mathfrak P^{\mathrm{pred}}_{n,m}(r,\mu,K,q,B,\lambda)}
 E_P[\mathcal L(\widetilde U,U)]
 \le
 \sup_{P\in\mathfrak P^{\mathrm{pred}}_{n,m}(r,\mu,K,q,B,\lambda)}
 E_P[\mathcal L(\widehat U^{\mathrm I},U)]
 \le C a^{\mathrm{loose}}_{nm},
\]
where
\[
 a^{\mathrm{loose}}_{nm}=
 \sqrt{\frac n{\mu r}}\left[1\wedge
 \frac B{\lambda_{nm}}\left\{
 \sqrt{\frac{(n\vee m)\log(e(n+m))}{q}}+
 \frac{\log(e(n+m))}{q}\right\}\right].
\]
The two endpoints use the same expected normalized maximum-row loss.  The bounds need not match: the endogenous-target obstruction rules out \(b_{nm}\)-attainment by recorded IPW--SVD, but does not rule it out for every estimator in \(\mathfrak E_{n,m}\).
%%% FIELD minimax-bracket-proof
The first inequality is exactly Theorem \(\mathrm{thm:adaptive\mbox{-}minimax\mbox{-}lower}\), whose lower-bound laws form a constant-propensity subclass of the same predictable class.  By \(\mathrm{def:observed\mbox{-}estimators}\), every member of \(\mathfrak E_{n,m}\) is measurable with respect to \(\mathcal O_{n,m}=\sigma(A,Y,p)\).  The recorded score is a measurable function of \((A,Y,p)\), and the deterministic rank projection and fixed tie rule are measurable; hence \(\widehat U^{\mathrm I}\in\mathfrak E_{n,m}\).  Evaluating the infimum at this admissible estimator proves the middle inequality.  Proposition \(\mathrm{prop:ipw\mbox{-}expected\mbox{-}risk\mbox{-}upper}\) proves the last inequality under the same class, loss, and sampling law.  Proposition \(\mathrm{prop:adaptive\mbox{-}endogenous\mbox{-}target\mbox{-}obstruction}\) concerns only the risk of \(\widehat U^{\mathrm I}\), so it cannot strengthen the first inequality to a lower bound above order \(b_{nm}\) for all estimators.  This proves the bracket and its stated scope.
%%% FIELD blind-resolution-statement
Uniformly over sequences \(P\in\mathfrak P^{\mathrm{pred}}_{n,m}(r,\mu,K,q,B,\lambda)\) that satisfy propensity-path stability, the primitive oracle blind gap, and \(b_{nm}=o(1)\), define
\[
 d_{nm}=\sqrt{\frac n{\mu r}}\left[1\wedge
 \frac B{\lambda_{nm}}\left\{
 q^{-1}\sqrt{(n\vee m)\log(e(n+m))}+q^{-3}\log(en)
 \right\}\right].
\]
Then
\[
 \left|\mathcal L(\widehat U^{\mathrm B},U)-\rho_{nm}\right|=O_P(d_{nm}),
 \qquad
 |\widehat\rho_{nm}-\rho_{nm}|=O_P(d_{nm}).
\]
Thus the observed blind-versus-recorded discrepancy is certified at the explicit \(d_{nm}\) resolution; no \(b_{nm}\)-scale equivalence is asserted.
%%% FIELD blind-resolution-proof
For \(Q_1,Q_2\in\mathbb O_{n,r}\), set
\[
 \delta(Q_1,Q_2)=\min_{R\in\mathbb O_r}
 \lVert Q_1R-Q_2\rVert_{2,\infty}.
\]
The minimum exists because \(\mathbb O_r\) is compact.  Right multiplication by an orthogonal matrix preserves every row Euclidean norm.  Consequently \(\delta\) is symmetric.  If \(R_{12}\) and \(R_{23}\) are minimizers for \((Q_1,Q_2)\) and \((Q_2,Q_3)\), respectively, then
\[
 \begin{aligned}
 \delta(Q_1,Q_3)
 &\le\lVert Q_1R_{12}R_{23}-Q_3\rVert_{2,\infty}\\
 &\le\lVert(Q_1R_{12}-Q_2)R_{23}\rVert_{2,\infty}
 +\lVert Q_2R_{23}-Q_3\rVert_{2,\infty}\\
 &=\delta(Q_1,Q_2)+\delta(Q_2,Q_3).
 \end{aligned}
\]
Multiplication by the positive constant \(\sqrt{n/(\mu r)}\) proves the triangle and reverse-triangle inequalities for \(\mathcal L\).

Lemma \(\mathrm{lem:blind\mbox{-}ratio\mbox{-}expansion}\), under exactly the hypotheses in the statement, gives
\[
 \mathcal L(\widehat U^{\mathrm B},U^{\mathrm B})=O_P(d_{nm}).
\]
Let \(a^{\mathrm{loose}}_{nm}\) denote the radius in Theorem \(\mathrm{thm:ipw\mbox{-}rowwise\mbox{-}upper}\).  Since \(q<1/2\),
\(q^{-1/2}\le q^{-1}\).  Bounded aspect ratio also gives a constant \(C\), depending only on \(c_-\) and \(c_+\), for which
\(\log(e(n+m))\le C\log(en)\).  Hence the expression inside the truncation defining \(a^{\mathrm{loose}}_{nm}\) is at most a constant times the one defining \(d_{nm}\).  Because
\(1\wedge(Cx)\le(1\vee C)(1\wedge x)\) for \(x\ge0\),
\[
 a^{\mathrm{loose}}_{nm}=O(d_{nm}).
\]
The recorded-IPW theorem therefore implies
\[
 \mathcal L(\widehat U^{\mathrm I},U)=O_P(d_{nm}).
\]

By the reverse-triangle inequality and the definitions of \(\rho_{nm}\) and \(\widehat\rho_{nm}\),
\[
 \begin{aligned}
 \left|\mathcal L(\widehat U^{\mathrm B},U)-\rho_{nm}\right|
 &\le\mathcal L(\widehat U^{\mathrm B},U^{\mathrm B}),\\
 |\widehat\rho_{nm}-\rho_{nm}|
 &=\left|\mathcal L(\widehat U^{\mathrm B},\widehat U^{\mathrm I})
 -\mathcal L(U^{\mathrm B},U)\right|\\
 &\le\mathcal L(\widehat U^{\mathrm B},U^{\mathrm B})
 +\mathcal L(\widehat U^{\mathrm I},U).
 \end{aligned}
\]
Substitution of the two stochastic bounds proves both conclusions.  No comparison between \(d_{nm}\) and \(b_{nm}\) is used.
%%% FIELD fixed-resolution-gap-statement
Suppose \(n\to\infty\), \(c_-\le m/n\le c_+\), and \(r,\mu,K,q\) are fixed, and define
\[
 d_{nm}=\sqrt{\frac n{\mu r}}\left[1\wedge
 \frac B{\lambda_{nm}}\left\{
 q^{-1}\sqrt{(n\vee m)\log(e(n+m))}+q^{-3}\log(en)
 \right\}\right].
\]
If \(b_{nm}=o(1)\), then the truncations in both \(b_{nm}\) and \(d_{nm}\) are eventually inactive and
\[
 b_{nm}\asymp\frac{B\sqrt{n\log(en)}}{\lambda_{nm}},
 \qquad
 d_{nm}\asymp\sqrt{\frac n{\mu r}}\,b_{nm}.
\]
In particular \(d_{nm}/b_{nm}\asymp\sqrt n\to\infty\), so \(d_{nm}=O(b_{nm})\) is not a nonempty fixed-complexity subregime under \(b_{nm}=o(1)\).
%%% FIELD fixed-resolution-gap-proof
By the geometric-rate definition and bounded aspect ratio,
\[
 \sqrt{n^{-1}+nm^{-2}}\asymp n^{-1/2}
 \quad\text{and}\quad
 \sqrt{nm}\sqrt{n^{-1}+nm^{-2}}\asymp\sqrt n.
\]
Because \(r,\mu,K,q\) are fixed positive constants, the untruncated expression in the lower/reference-scale definition is therefore
\[
 \frac{B\sqrt{nm}}{\lambda_{nm}}\sqrt{\log(en)}\,b^0_{nm}
 \asymp\frac{B\sqrt{n\log(en)}}{\lambda_{nm}}.
\]
The hypothesis \(b_{nm}=o(1)\) makes this expression smaller than one eventually, proving the first equivalence and showing that the truncation defining \(b_{nm}\) is inactive.

Set
\[
 x_{nm}=\frac B{\lambda_{nm}}
 \left\{q^{-1}\sqrt{(n\vee m)\log(e(n+m))}
 +q^{-3}\log(en)\right\}.
\]
Bounded aspect ratio gives
\[
 \sqrt{(n\vee m)\log(e(n+m))}\asymp\sqrt{n\log(en)},
\]
whereas
\[
 \frac{\log(en)}{\sqrt{n\log(en)}}
 =\sqrt{\frac{\log(en)}n}\longrightarrow0.
\]
With fixed \(q\), it follows that
\[
 x_{nm}\asymp\frac{B\sqrt{n\log(en)}}{\lambda_{nm}}asymp b_{nm}=o(1).
\]
Thus the truncation in \(d_{nm}\) is also eventually inactive, and its definition gives
\[
 d_{nm}=\sqrt{\frac n{\mu r}}\,x_{nm}
 \asymp\sqrt{\frac n{\mu r}}\,b_{nm}.
\]
Since \(r\) and \(\mu\) are fixed and \(n\to\infty\), division by the positive \(b_{nm}\) proves the divergence and the final assertion.
